Edmond Halley, contemporani d'Isaac Newton, va descobrir que tres cometes descrits el 1531, el 1607 i el 1682 eren, en realitat, un mateix cometa, el cometa Halley. Aquest cometa descriu una òrbita el·líptica amb un període de 76 anys, aproximadament.

\begin{apartat}{1,25}
Determineu la longitud del semieix major de l'òrbita del cometa Halley. Representeu esquemàticament l'òrbita del cometa, indicant la posició del Sol, l'afeli (punt més allunyat) i el periheli (punt més proper).
\end{apartat}

\begin{apartat}{1,25}
El periheli es troba a 0,586 ua del centre del Sol. Determineu la intensitat de camp gravitatori a l'afeli i al periheli.
\end{apartat}

\dades{%
  $G = 6,67\times 10^{-11}\ \mathrm{N\,m^2\,kg^{-2}}$.\\
  $m_\mathrm{Sol} = 1,99\times 10^{30}\ \mathrm{kg}$.\\
  $1\ \mathrm{ua} = 1,50\times 10^{11}\ \mathrm{m}$.}
